\documentclass[11pt,a4paper]{article}
\usepackage[T1]{fontenc}
\usepackage[utf8]{inputenc}
\usepackage[french]{babel}
\usepackage{lmodern}
\usepackage{microtype}
\usepackage[margin=2.2cm]{geometry}
\usepackage{amsmath,amssymb,amsthm,mathtools}
\usepackage{icomma}
\usepackage{booktabs,tabularx,array}
\usepackage{xcolor}
\usepackage{tikz}
\usetikzlibrary{patterns,positioning}
\usepackage{pgfplots}
\pgfplotsset{compat=1.18}
\usepackage[most]{tcolorbox}
\usepackage{enumitem}
\usepackage[hidelinks]{hyperref}
\usepackage{caption}
\usepackage{float}
\usepackage{pdflscape}
\usepackage{pifont}

\definecolor{nyt}{HTML}{1F3B63}
\definecolor{wp}{HTML}{17695B}
\definecolor{gov}{HTML}{B42318}
\definecolor{nat}{HTML}{9A6212}
\definecolor{regnone}{HTML}{C6C2B8}
\definecolor{regnosuit}{HTML}{CBE2C3}
\definecolor{regnyt}{HTML}{F0C5BD}
\definecolor{regboth}{HTML}{F1D69B}

\newtcolorbox{cours}{colback=nat!6,colframe=nat!70,boxrule=0.5pt,arc=2pt,left=6pt,right=6pt,top=4pt,bottom=4pt,fonttitle=\bfseries\small,title=Lien avec le cours}
\newtcolorbox{resultat}{colback=wp!6,colframe=wp!70,boxrule=0.5pt,arc=2pt,left=6pt,right=6pt,top=4pt,bottom=4pt}
\newtcolorbox{intuition}{colback=nyt!4,colframe=nyt!50,boxrule=0.4pt,arc=2pt,left=6pt,right=6pt,top=3pt,bottom=3pt,fonttitle=\bfseries\small,title=Intuition}

\theoremstyle{definition}
\newtheorem{lemme}{Lemme}
\newtheorem{proposition}{Proposition}
\newcommand{\ind}[1]{\mathbf{1}_{\{\text{#1}\}}}
\newcommand{\E}{\mathbb{E}}
\setlength{\parindent}{0pt}
\setlength{\parskip}{5pt}
\captionsetup{font=small,labelfont=bf}

\newtcolorbox{erreur}{colback=gov!5,colframe=gov!70,boxrule=0.5pt,arc=2pt,left=6pt,right=6pt,top=4pt,bottom=4pt,fonttitle=\bfseries\small,title=Erreur}
\newtcolorbox{discuter}{colback=nat!6,colframe=nat!70,boxrule=0.5pt,arc=2pt,left=6pt,right=6pt,top=4pt,bottom=4pt,fonttitle=\bfseries\small,title=À discuter}
\newtcolorbox{simple}{colback=nyt!4,colframe=nyt!50,boxrule=0.4pt,arc=2pt,left=6pt,right=6pt,top=3pt,bottom=3pt,fonttitle=\bfseries\small,title=En mots}
\newcommand{\T}[1]{\textcolor{nyt}{#1}}
\newcommand{\W}[1]{\textcolor{wp}{#1}}
\newcommand{\N}[1]{\textcolor{gov}{#1}}


\title{\textbf{Pentagon Papers}\\[3pt]\large Scénario « bonus de Nixon » : le $+10$ rendu cohérent}
\author{Travail de théorie des jeux, groupe 5--6 · ECGEB366}
\date{}

\begin{document}
\maketitle
\vspace{-1.2cm}

\begin{resultat}
\textbf{En bref.}
\begin{itemize}[leftmargin=*,itemsep=1pt,topsep=2pt]
\item Le collègue donne $+10$ à Nixon quand personne ne publie ($O_7$), car c'est sa meilleure situation. L'idée est juste.
\item Pour être cohérent, ce $+10$ doit devenir un \textbf{bonus $R=10$ dans l'équation de Nixon}, appliqué à \emph{toutes} les feuilles où le secret est gardé.
\item Équation retenue : $u_N=R\,(1-\sigma)-D\,\sigma-k_T\,\ind{Times poursuivi}-k_W\,\ind{Post poursuivi}$.
\item Avec $D=200$, le bonus \textbf{ne change pas l'équilibre} : $O_7$ sans $\theta$, $O_1$ (1971) avec $\theta=0,3$.
\item Il compte quand le dommage $D$ est petit : avec $D=5$ et $\theta=0,3$, il fait passer de $O_4$ (aucun procès) à $O_1$ (1971).
\end{itemize}
\end{resultat}

\begin{discuter}
\textbf{Lettres.} Comme dans le dessin du collègue : $p$ = chance que le \textbf{juge local} refuse la censure, $q$ = chance que la \textbf{Cour suprême} donne raison aux journaux. Exemple chiffré : $p=0,3$, $q=0,65$.
\end{discuter}

\section{Le problème de départ}
Dans le dessin, Nixon a $+10$ en $O_7$ (personne ne publie), mais $-1$ quand une censure tient et que rien ne sort non plus. Deux situations où le secret est gardé reçoivent donc des valeurs très différentes, sans raison. Le $+10$ n'est pas une erreur d'idée : c'est un \textbf{terme qui manque dans l'équation}.

\begin{simple}
Règle d'or : une seule équation par joueur, appliquée à toutes les feuilles. Si Nixon gagne quelque chose quand le secret est gardé, il le gagne \emph{chaque fois} que le secret est gardé.
\end{simple}

\section{L'équation de Nixon avec bonus}

\subsection{La forme exacte}
\[
u_N=\underbrace{R\,(1-\sigma)}_{\text{bonus : part gardée secrète}}\;-\;\underbrace{D\,\sigma}_{\text{dommage : part qui sort}}
\]
\[
\hphantom{u_N={}}\;-\;\underbrace{k_T\,\ind{Times poursuivi}}_{\text{procès du Times}}\;-\;\underbrace{k_W\,\ind{Post poursuivi}}_{\text{procès du Post}}
\]
avec $\sigma$ la part des documents qui sort :
\begin{center}
\small
\begin{tabular}{@{}llll@{}}
\toprule
Situation & $\sigma$ & Part gardée secrète $1-\sigma$ & Bonus $R(1-\sigma)$ \\
\midrule
Tout sort & $1$ & $0$ & $0$ \\
Censure maintenue & $\theta$ & $1-\theta$ & $R(1-\theta)$ \\
Rien ne sort & $0$ & $1$ & $R$ \\
\bottomrule
\end{tabular}
\end{center}

\subsection{Les deux cas}
\begin{itemize}[leftmargin=*,itemsep=2pt]
\item \textbf{Sans $\theta$} ($\theta=0$) : une censure veut dire que \emph{rien ne sort}. Le bonus est entier, $R=10$, dans $O_7$ \emph{et} après chaque censure. L'équation devient $u_N=R\,\ind{rien ne sort}-D\,\ind{tout sort}-k_T\,\ind{\dots}-k_W\,\ind{\dots}$.
\item \textbf{Avec $\theta=0,3$} : une censure laisse sortir $30\,\%$. Nixon reçoit le bonus pour les $70\,\%$ gardés secrets : $R(1-\theta)=7$. Seul $O_7$ donne le bonus entier.
\end{itemize}
Les gains des journaux ne changent pas : le bonus ne concerne que Nixon.

\subsection{Les valeurs}
\[
B_T=B_W=10,\quad c_T=1,\quad c_W=5,\quad S_T=S_W=5,\quad D=200,\quad k_T=k_W=1,
\]
\[
\mathbf{R=10},\qquad p=0,3,\qquad q=0,65 .
\]

\subsection{Les gains de Nixon après chaque situation}
\begin{center}
\small
\begin{tabular}{@{}llll@{}}
\toprule
Situation & Formule & Sans $\theta$ & Avec $\theta=0,3$ \\
\midrule
Rien ne sort, aucun procès & $R$ & $10$ & $10$ \\
Censure, 1 procès & $R(1-\theta)-\theta D-k$ & $10-1=\mathbf{9}$ & $7-60-1=\mathbf{-54}$ \\
Censure, 2 procès & $R(1-\theta)-\theta D-k_T-k_W$ & $10-2=\mathbf{8}$ & $7-60-2=\mathbf{-55}$ \\
Tout sort, aucun procès & $-D$ & $-200$ & $-200$ \\
Tout sort, 1 procès & $-D-k$ & $-201$ & $-201$ \\
Tout sort, 2 procès & $-D-k_T-k_W$ & $-202$ & $-202$ \\
\bottomrule
\end{tabular}
\end{center}
Pour comparer, sans bonus : censure à $-1$ et $-2$ (sans $\theta$), à $-61$ et $-62$ (avec $\theta$).

\subsection{Les gains des journaux (inchangés)}
\begin{center}
\small
\begin{tabular}{@{}llll@{}}
\toprule
Situation & Formule & Sans $\theta$ & Avec $\theta=0,3$ \\
\midrule
\T{Times gagne} & $B_T-c_T$ & $9$ & $9$ \\
\T{Times perd} & $\theta B_T-c_T-S_T$ & $-6$ & $-3$ \\
\W{Post gagne} & $B_W-c_W$ & $5$ & $5$ \\
\W{Post perd} & $\theta B_W-c_W-S_W$ & $-10$ & $-7$ \\
\bottomrule
\end{tabular}
\end{center}

\begin{landscape}
\begin{center}{\large\bfseries L'arbre avec bonus, en lettres}\\[2pt]{\small Gains dans l'ordre \textcolor{nyt}{\textbf{Times}}, \textcolor{wp}{\textbf{Post}}, \textcolor{gov}{\textbf{Nixon}}. $p$ = juge local, $q$ = Cour suprême. En gras : là où le bonus $R$ intervient. Sans $\theta$, poser $\theta=0$ : le terme $R(1-\theta)$ devient $R$.}\end{center}\vspace{-4pt}
\noindent\resizebox{\linewidth}{!}{\begin{tikzpicture}[x=1cm,y=1cm,lab/.style={font=\scriptsize,fill=white,inner sep=1pt,sloped},eq/.style={line width=1.8pt,draw=gov},br/.style={line width=0.6pt,draw=black!55}]
\draw[br] (0.4,-7.928125)--(2.6,-4.56875) node[lab,midway,above]{Publier};
\draw[br] (0.4,-7.928125)--(5.0,-11.2875) node[lab,midway,below]{Ne pas publier};
\draw[br] (2.6,-4.56875)--(5.0,-2.6875) node[lab,midway,above]{Poursuivre};
\draw[br] (2.6,-4.56875)--(9.6,-6.449999999999999) node[lab,midway,below]{Laisser faire};
\draw[br] (5.0,-2.6875)--(7.4,-1.075) node[lab,midway,above]{Publier};
\draw[br] (5.0,-2.6875)--(9.6,-4.3) node[lab,midway,below]{Ne pas publier};
\draw[br] (7.4,-1.075)--(9.6,0.0) node[lab,midway,above]{Poursuivre};
\draw[br] (7.4,-1.075)--(9.6,-2.15) node[lab,midway,below]{Laisser faire};
\draw[br] (5.0,-11.2875)--(7.4,-9.675) node[lab,midway,above]{Publier};
\draw[br] (5.0,-11.2875)--(9.6,-12.899999999999999) node[lab,midway,below]{Ne pas publier};
\draw[br] (7.4,-9.675)--(9.6,-8.6) node[lab,midway,above]{Poursuivre};
\draw[br] (7.4,-9.675)--(9.6,-10.75) node[lab,midway,below]{Laisser faire};
\fill[nyt] (0.4,-7.928125) circle (4pt); \node[font=\small\bfseries,text=nyt,above left=3pt and -2pt] at (0.4,-7.928125) {Times};
\fill[gov] (2.6,-4.56875) circle (4pt); \node[font=\small\bfseries,text=gov,above left=3pt and -2pt] at (2.6,-4.56875) {Nixon};
\fill[wp] (5.0,-2.6875) circle (4pt); \node[font=\small\bfseries,text=wp,above left=3pt and -2pt] at (5.0,-2.6875) {Post};
\fill[gov] (7.4,-1.075) circle (4pt); \node[font=\small\bfseries,text=gov,above left=3pt and -2pt] at (7.4,-1.075) {Nixon};
\fill[wp] (5.0,-11.2875) circle (4pt); \node[font=\small\bfseries,text=wp,above left=3pt and -2pt] at (5.0,-11.2875) {Post};
\fill[gov] (7.4,-9.675) circle (4pt); \node[font=\small\bfseries,text=gov,above left=3pt and -2pt] at (7.4,-9.675) {Nixon};
\fill (9.6,0.0) circle (2pt);
\node[anchor=west,draw=black!35,fill=white,rounded corners=3pt,text width=15.2cm,align=left,font=\footnotesize,inner sep=4pt] at (9.85,0.0) {\textbf{$O_1$ · les deux journaux poursuivis}\\\textcolor{nyt}{\textbf{Times} : $p\,(B_T-c_T)+(1-p)\,q\,(B_T-c_T)+(1-p)(1-q)\,(\theta B_T-c_T-S_T)$}\\\textcolor{wp}{\textbf{Post} : $p\cdot0+(1-p)\,q\,(B_W-c_W)+(1-p)(1-q)\,(\theta B_W-c_W-S_W)$}\\\textcolor{gov}{\textbf{Nixon} : $p\,(-D-k_T)+(1-p)\,q\,(-D-k_T-k_W)+(1-p)(1-q)\,\big(\mathbf{R(1-\theta)}-\theta D-k_T-k_W\big)$}};
\fill (9.6,-2.15) circle (2pt);
\node[anchor=west,draw=black!35,fill=white,rounded corners=3pt,text width=15.2cm,align=left,font=\footnotesize,inner sep=4pt] at (9.85,-2.15) {\textbf{$O_2$ · Times poursuivi, Post libre}\\\textcolor{nyt}{\textbf{Times} : $B_T-c_T$}\\\textcolor{wp}{\textbf{Post} : $p\cdot0+(1-p)\,B_W$}\\\textcolor{gov}{\textbf{Nixon} : $-D-k_T$}};
\fill (9.6,-4.3) circle (2pt);
\node[anchor=west,draw=black!35,fill=white,rounded corners=3pt,text width=15.2cm,align=left,font=\footnotesize,inner sep=4pt] at (9.85,-4.3) {\textbf{$O_3$ · Times poursuivi, Post muet}\\\textcolor{nyt}{\textbf{Times} : $p\,(B_T-c_T)+(1-p)\,(\theta B_T-c_T-S_T)$}\\\textcolor{wp}{\textbf{Post} : $0$}\\\textcolor{gov}{\textbf{Nixon} : $p\,(-D-k_T)+(1-p)\,\big(\mathbf{R(1-\theta)}-\theta D-k_T\big)$}};
\fill (9.6,-6.449999999999999) circle (2pt);
\node[anchor=west,draw=black!35,fill=white,rounded corners=3pt,text width=15.2cm,align=left,font=\footnotesize,inner sep=4pt] at (9.85,-6.449999999999999) {\textbf{$O_4$ · Times publie, pas de procès}\\\textcolor{nyt}{\textbf{Times} : $B_T$}\\\textcolor{wp}{\textbf{Post} : $0$}\\\textcolor{gov}{\textbf{Nixon} : $-D$}};
\fill (9.6,-8.6) circle (2pt);
\node[anchor=west,draw=black!35,fill=white,rounded corners=3pt,text width=15.2cm,align=left,font=\footnotesize,inner sep=4pt] at (9.85,-8.6) {\textbf{$O_5$ · Post seul, poursuivi}\\\textcolor{nyt}{\textbf{Times} : $0$}\\\textcolor{wp}{\textbf{Post} : $p\,(B_W-c_W)+(1-p)\,(\theta B_W-c_W-S_W)$}\\\textcolor{gov}{\textbf{Nixon} : $p\,(-D-k_W)+(1-p)\,\big(\mathbf{R(1-\theta)}-\theta D-k_W\big)$}};
\fill (9.6,-10.75) circle (2pt);
\node[anchor=west,draw=black!35,fill=white,rounded corners=3pt,text width=15.2cm,align=left,font=\footnotesize,inner sep=4pt] at (9.85,-10.75) {\textbf{$O_6$ · Post seul, pas de procès}\\\textcolor{nyt}{\textbf{Times} : $0$}\\\textcolor{wp}{\textbf{Post} : $B_W$}\\\textcolor{gov}{\textbf{Nixon} : $-D$}};
\fill (9.6,-12.899999999999999) circle (2pt);
\node[anchor=west,draw=black!35,fill=white,rounded corners=3pt,text width=15.2cm,align=left,font=\footnotesize,inner sep=4pt] at (9.85,-12.899999999999999) {\textbf{$O_7$ · personne ne publie}\\\textcolor{nyt}{\textbf{Times} : $0$}\\\textcolor{wp}{\textbf{Post} : $0$}\\\textcolor{gov}{\textbf{Nixon} : $\mathbf{R}$}};
\end{tikzpicture}}
\end{landscape}

\section{Scénario A : bonus, sans \texorpdfstring{$\theta$}{theta}}

\begin{landscape}
\begin{center}{\large\bfseries Scénario A : bonus $R=10$, sans $\theta$}\\[2pt]{\small Gains dans l'ordre \textcolor{nyt}{\textbf{Times}}, \textcolor{wp}{\textbf{Post}}, \textcolor{gov}{\textbf{Nixon}}. $p$ = juge local, $q$ = Cour suprême. En gras : là où le bonus $R$ intervient. Traits rouges épais : choix à l'équilibre $\Rightarrow$ $O_7$.}\end{center}\vspace{-4pt}
\noindent\resizebox{\linewidth}{!}{\begin{tikzpicture}[x=1cm,y=1cm,lab/.style={font=\scriptsize,fill=white,inner sep=1pt,sloped},eq/.style={line width=1.8pt,draw=gov},br/.style={line width=0.6pt,draw=black!55}]
\draw[br] (0.4,-7.928125)--(2.6,-4.56875) node[lab,midway,above]{Publier};
\draw[eq] (0.4,-7.928125)--(5.0,-11.2875) node[lab,midway,below]{Ne pas publier};
\draw[eq] (2.6,-4.56875)--(5.0,-2.6875) node[lab,midway,above]{Poursuivre};
\draw[br] (2.6,-4.56875)--(9.6,-6.449999999999999) node[lab,midway,below]{Laisser faire};
\draw[br] (5.0,-2.6875)--(7.4,-1.075) node[lab,midway,above]{Publier};
\draw[eq] (5.0,-2.6875)--(9.6,-4.3) node[lab,midway,below]{Ne pas publier};
\draw[eq] (7.4,-1.075)--(9.6,0.0) node[lab,midway,above]{Poursuivre};
\draw[br] (7.4,-1.075)--(9.6,-2.15) node[lab,midway,below]{Laisser faire};
\draw[br] (5.0,-11.2875)--(7.4,-9.675) node[lab,midway,above]{Publier};
\draw[eq] (5.0,-11.2875)--(9.6,-12.899999999999999) node[lab,midway,below]{Ne pas publier};
\draw[eq] (7.4,-9.675)--(9.6,-8.6) node[lab,midway,above]{Poursuivre};
\draw[br] (7.4,-9.675)--(9.6,-10.75) node[lab,midway,below]{Laisser faire};
\fill[nyt] (0.4,-7.928125) circle (4pt); \node[font=\small\bfseries,text=nyt,above left=3pt and -2pt] at (0.4,-7.928125) {Times};
\fill[gov] (2.6,-4.56875) circle (4pt); \node[font=\small\bfseries,text=gov,above left=3pt and -2pt] at (2.6,-4.56875) {Nixon};
\fill[wp] (5.0,-2.6875) circle (4pt); \node[font=\small\bfseries,text=wp,above left=3pt and -2pt] at (5.0,-2.6875) {Post};
\fill[gov] (7.4,-1.075) circle (4pt); \node[font=\small\bfseries,text=gov,above left=3pt and -2pt] at (7.4,-1.075) {Nixon};
\fill[wp] (5.0,-11.2875) circle (4pt); \node[font=\small\bfseries,text=wp,above left=3pt and -2pt] at (5.0,-11.2875) {Post};
\fill[gov] (7.4,-9.675) circle (4pt); \node[font=\small\bfseries,text=gov,above left=3pt and -2pt] at (7.4,-9.675) {Nixon};
\fill (9.6,0.0) circle (2pt);
\node[anchor=west,draw=black!35,fill=white,rounded corners=3pt,text width=15.2cm,align=left,font=\footnotesize,inner sep=4pt] at (9.85,0.0) {\textbf{$O_1$ · les deux journaux poursuivis}\\\textcolor{nyt}{\textbf{Times} : $p\cdot9+(1-p)\,q\cdot9+(1-p)(1-q)\cdot(-6)$ $=\mathbf{5,325}$}\\\textcolor{wp}{\textbf{Post} : $p\cdot0+(1-p)\,q\cdot5+(1-p)(1-q)\cdot(-10)$ $=\mathbf{-0,175}$}\\\textcolor{gov}{\textbf{Nixon} : $p\cdot(-201)+(1-p)\,q\cdot(-202)+(1-p)(1-q)\cdot\mathbf{8}$ $=\mathbf{-150,25}$}};
\fill (9.6,-2.15) circle (2pt);
\node[anchor=west,draw=black!35,fill=white,rounded corners=3pt,text width=15.2cm,align=left,font=\footnotesize,inner sep=4pt] at (9.85,-2.15) {\textbf{$O_2$ · Times poursuivi, Post libre}\\\textcolor{nyt}{\textbf{Times} : $9$}\\\textcolor{wp}{\textbf{Post} : $p\cdot0+(1-p)\cdot10$ $=\mathbf{7}$}\\\textcolor{gov}{\textbf{Nixon} : $-201$}};
\fill (9.6,-4.3) circle (2pt);
\node[anchor=west,draw=black!35,fill=white,rounded corners=3pt,text width=15.2cm,align=left,font=\footnotesize,inner sep=4pt] at (9.85,-4.3) {\textbf{$O_3$ · Times poursuivi, Post muet}\\\textcolor{nyt}{\textbf{Times} : $p\cdot9+(1-p)\cdot(-6)$ $=\mathbf{-1,5}$}\\\textcolor{wp}{\textbf{Post} : $0$}\\\textcolor{gov}{\textbf{Nixon} : $p\cdot(-201)+(1-p)\cdot\mathbf{9}$ $=\mathbf{-54}$}};
\fill (9.6,-6.449999999999999) circle (2pt);
\node[anchor=west,draw=black!35,fill=white,rounded corners=3pt,text width=15.2cm,align=left,font=\footnotesize,inner sep=4pt] at (9.85,-6.449999999999999) {\textbf{$O_4$ · Times publie, pas de procès}\\\textcolor{nyt}{\textbf{Times} : $10$}\\\textcolor{wp}{\textbf{Post} : $0$}\\\textcolor{gov}{\textbf{Nixon} : $-200$}};
\fill (9.6,-8.6) circle (2pt);
\node[anchor=west,draw=black!35,fill=white,rounded corners=3pt,text width=15.2cm,align=left,font=\footnotesize,inner sep=4pt] at (9.85,-8.6) {\textbf{$O_5$ · Post seul, poursuivi}\\\textcolor{nyt}{\textbf{Times} : $0$}\\\textcolor{wp}{\textbf{Post} : $p\cdot5+(1-p)\cdot(-10)$ $=\mathbf{-5,5}$}\\\textcolor{gov}{\textbf{Nixon} : $p\cdot(-201)+(1-p)\cdot\mathbf{9}$ $=\mathbf{-54}$}};
\fill (9.6,-10.75) circle (2pt);
\node[anchor=west,draw=black!35,fill=white,rounded corners=3pt,text width=15.2cm,align=left,font=\footnotesize,inner sep=4pt] at (9.85,-10.75) {\textbf{$O_6$ · Post seul, pas de procès}\\\textcolor{nyt}{\textbf{Times} : $0$}\\\textcolor{wp}{\textbf{Post} : $10$}\\\textcolor{gov}{\textbf{Nixon} : $-200$}};
\fill (9.6,-12.899999999999999) circle (2pt);
\node[anchor=west,draw=black!35,fill=white,rounded corners=3pt,text width=15.2cm,align=left,font=\footnotesize,inner sep=4pt] at (9.85,-12.899999999999999) {\textbf{$O_7$ · personne ne publie}\\\textcolor{nyt}{\textbf{Times} : $0$}\\\textcolor{wp}{\textbf{Post} : $0$}\\\textcolor{gov}{\textbf{Nixon} : $\mathbf{10}$}};
\end{tikzpicture}}
\end{landscape}

\subsection{Les feuilles de Nixon qui changent}
\begin{align*}
O_1:\quad \N{\E[u_N]} &= p\,(-D-k_T)+(1-p)\,q\,(-D-k_T-k_W)+(1-p)(1-q)\,(R-k_T-k_W)\\
&= 0,3\times(-201)+0,455\times(-202)+0,245\times8=-60,3-91,91+1,96=\mathbf{-150,25}\\[4pt]
O_3:\quad \N{\E[u_N]} &= p\,(-D-k_T)+(1-p)\,(R-k_T)=0,3\times(-201)+0,7\times9=-60,3+6,3=\mathbf{-54}\\[4pt]
O_5:\quad \N{\E[u_N]} &= p\,(-D-k_W)+(1-p)\,(R-k_W)=\mathbf{-54}\\[4pt]
O_7:\quad \N{u_N} &= R=\mathbf{10}
\end{align*}
Sans bonus, ces valeurs étaient $-152,7$ ; $-61$ ; $-61$ ; $0$.

\subsection{La résolution}
\begin{center}
\small
\begin{tabular}{@{}lll@{}}
\toprule
Qui décide & Comparaison & Choix \\
\midrule
\N{Nixon} face au Post en relais & $O_1:-150,25$ contre $O_2:-201$ & poursuit \\
\W{Post} en relais & $O_1:-0,175$ contre $O_3:0$ & se tait \\
\N{Nixon} face au Times & $O_3:-54$ contre $O_4:-200$ & poursuit \\
\N{Nixon} face au Post seul & $O_5:-54$ contre $O_6:-200$ & poursuit \\
\W{Post} seul & $O_5:-5,5$ contre $O_7:0$ & se tait \\
\T{Times} & $O_3:-1,5$ contre $O_7:0$ & ne publie pas \\
\bottomrule
\end{tabular}
\end{center}
\begin{resultat}
\textbf{Résultat A : $O_7$, rien ne sort.} Même résultat que sans bonus : Nixon poursuivait déjà, et les journaux décident avec leurs propres gains, qui n'ont pas changé.
\end{resultat}

\section{Scénario B : bonus, avec \texorpdfstring{$\theta=0,3$}{theta = 0,3}}

\begin{landscape}
\begin{center}{\large\bfseries Scénario B : bonus $R=10$, avec $\theta=0,3$}\\[2pt]{\small Gains dans l'ordre \textcolor{nyt}{\textbf{Times}}, \textcolor{wp}{\textbf{Post}}, \textcolor{gov}{\textbf{Nixon}}. $p$ = juge local, $q$ = Cour suprême. En gras : là où le bonus $R$ intervient. Traits rouges épais : choix à l'équilibre $\Rightarrow$ $O_1$, c'est 1971.}\end{center}\vspace{-4pt}
\noindent\resizebox{\linewidth}{!}{\begin{tikzpicture}[x=1cm,y=1cm,lab/.style={font=\scriptsize,fill=white,inner sep=1pt,sloped},eq/.style={line width=1.8pt,draw=gov},br/.style={line width=0.6pt,draw=black!55}]
\draw[eq] (0.4,-7.928125)--(2.6,-4.56875) node[lab,midway,above]{Publier};
\draw[br] (0.4,-7.928125)--(5.0,-11.2875) node[lab,midway,below]{Ne pas publier};
\draw[eq] (2.6,-4.56875)--(5.0,-2.6875) node[lab,midway,above]{Poursuivre};
\draw[br] (2.6,-4.56875)--(9.6,-6.449999999999999) node[lab,midway,below]{Laisser faire};
\draw[eq] (5.0,-2.6875)--(7.4,-1.075) node[lab,midway,above]{Publier};
\draw[br] (5.0,-2.6875)--(9.6,-4.3) node[lab,midway,below]{Ne pas publier};
\draw[eq] (7.4,-1.075)--(9.6,0.0) node[lab,midway,above]{Poursuivre};
\draw[br] (7.4,-1.075)--(9.6,-2.15) node[lab,midway,below]{Laisser faire};
\draw[br] (5.0,-11.2875)--(7.4,-9.675) node[lab,midway,above]{Publier};
\draw[eq] (5.0,-11.2875)--(9.6,-12.899999999999999) node[lab,midway,below]{Ne pas publier};
\draw[eq] (7.4,-9.675)--(9.6,-8.6) node[lab,midway,above]{Poursuivre};
\draw[br] (7.4,-9.675)--(9.6,-10.75) node[lab,midway,below]{Laisser faire};
\fill[nyt] (0.4,-7.928125) circle (4pt); \node[font=\small\bfseries,text=nyt,above left=3pt and -2pt] at (0.4,-7.928125) {Times};
\fill[gov] (2.6,-4.56875) circle (4pt); \node[font=\small\bfseries,text=gov,above left=3pt and -2pt] at (2.6,-4.56875) {Nixon};
\fill[wp] (5.0,-2.6875) circle (4pt); \node[font=\small\bfseries,text=wp,above left=3pt and -2pt] at (5.0,-2.6875) {Post};
\fill[gov] (7.4,-1.075) circle (4pt); \node[font=\small\bfseries,text=gov,above left=3pt and -2pt] at (7.4,-1.075) {Nixon};
\fill[wp] (5.0,-11.2875) circle (4pt); \node[font=\small\bfseries,text=wp,above left=3pt and -2pt] at (5.0,-11.2875) {Post};
\fill[gov] (7.4,-9.675) circle (4pt); \node[font=\small\bfseries,text=gov,above left=3pt and -2pt] at (7.4,-9.675) {Nixon};
\fill (9.6,0.0) circle (2pt);
\node[anchor=west,draw=black!35,fill=white,rounded corners=3pt,text width=15.2cm,align=left,font=\footnotesize,inner sep=4pt] at (9.85,0.0) {\textbf{$O_1$ · les deux journaux poursuivis}\\\textcolor{nyt}{\textbf{Times} : $p\cdot9+(1-p)\,q\cdot9+(1-p)(1-q)\cdot(-3)$ $=\mathbf{6,06}$}\\\textcolor{wp}{\textbf{Post} : $p\cdot0+(1-p)\,q\cdot5+(1-p)(1-q)\cdot(-7)$ $=\mathbf{0,56}$}\\\textcolor{gov}{\textbf{Nixon} : $p\cdot(-201)+(1-p)\,q\cdot(-202)+(1-p)(1-q)\cdot(\mathbf{-55})$ $=\mathbf{-165,69}$}};
\fill (9.6,-2.15) circle (2pt);
\node[anchor=west,draw=black!35,fill=white,rounded corners=3pt,text width=15.2cm,align=left,font=\footnotesize,inner sep=4pt] at (9.85,-2.15) {\textbf{$O_2$ · Times poursuivi, Post libre}\\\textcolor{nyt}{\textbf{Times} : $9$}\\\textcolor{wp}{\textbf{Post} : $p\cdot0+(1-p)\cdot10$ $=\mathbf{7}$}\\\textcolor{gov}{\textbf{Nixon} : $-201$}};
\fill (9.6,-4.3) circle (2pt);
\node[anchor=west,draw=black!35,fill=white,rounded corners=3pt,text width=15.2cm,align=left,font=\footnotesize,inner sep=4pt] at (9.85,-4.3) {\textbf{$O_3$ · Times poursuivi, Post muet}\\\textcolor{nyt}{\textbf{Times} : $p\cdot9+(1-p)\cdot(-3)$ $=\mathbf{0,6}$}\\\textcolor{wp}{\textbf{Post} : $0$}\\\textcolor{gov}{\textbf{Nixon} : $p\cdot(-201)+(1-p)\cdot(\mathbf{-54})$ $=\mathbf{-98,1}$}};
\fill (9.6,-6.449999999999999) circle (2pt);
\node[anchor=west,draw=black!35,fill=white,rounded corners=3pt,text width=15.2cm,align=left,font=\footnotesize,inner sep=4pt] at (9.85,-6.449999999999999) {\textbf{$O_4$ · Times publie, pas de procès}\\\textcolor{nyt}{\textbf{Times} : $10$}\\\textcolor{wp}{\textbf{Post} : $0$}\\\textcolor{gov}{\textbf{Nixon} : $-200$}};
\fill (9.6,-8.6) circle (2pt);
\node[anchor=west,draw=black!35,fill=white,rounded corners=3pt,text width=15.2cm,align=left,font=\footnotesize,inner sep=4pt] at (9.85,-8.6) {\textbf{$O_5$ · Post seul, poursuivi}\\\textcolor{nyt}{\textbf{Times} : $0$}\\\textcolor{wp}{\textbf{Post} : $p\cdot5+(1-p)\cdot(-7)$ $=\mathbf{-3,4}$}\\\textcolor{gov}{\textbf{Nixon} : $p\cdot(-201)+(1-p)\cdot(\mathbf{-54})$ $=\mathbf{-98,1}$}};
\fill (9.6,-10.75) circle (2pt);
\node[anchor=west,draw=black!35,fill=white,rounded corners=3pt,text width=15.2cm,align=left,font=\footnotesize,inner sep=4pt] at (9.85,-10.75) {\textbf{$O_6$ · Post seul, pas de procès}\\\textcolor{nyt}{\textbf{Times} : $0$}\\\textcolor{wp}{\textbf{Post} : $10$}\\\textcolor{gov}{\textbf{Nixon} : $-200$}};
\fill (9.6,-12.899999999999999) circle (2pt);
\node[anchor=west,draw=black!35,fill=white,rounded corners=3pt,text width=15.2cm,align=left,font=\footnotesize,inner sep=4pt] at (9.85,-12.899999999999999) {\textbf{$O_7$ · personne ne publie}\\\textcolor{nyt}{\textbf{Times} : $0$}\\\textcolor{wp}{\textbf{Post} : $0$}\\\textcolor{gov}{\textbf{Nixon} : $\mathbf{10}$}};
\end{tikzpicture}}
\end{landscape}

\subsection{Les feuilles de Nixon qui changent}
\begin{align*}
O_1:\quad \N{\E[u_N]} &= p\,(-D-k_T)+(1-p)\,q\,(-D-k_T-k_W)\\
&\quad+(1-p)(1-q)\,\big(R(1-\theta)-\theta D-k_T-k_W\big)\\
&= 0,3\times(-201)+0,455\times(-202)+0,245\times(-55)=-60,3-91,91-13,475=\mathbf{-165,69}\\[4pt]
O_3:\quad \N{\E[u_N]} &= p\,(-D-k_T)+(1-p)\,\big(R(1-\theta)-\theta D-k_T\big)\\
&=0,3\times(-201)+0,7\times(-54)=-60,3-37,8=\mathbf{-98,1}\\[4pt]
O_5:\quad \N{\E[u_N]} &= p\,(-D-k_W)+(1-p)\,\big(R(1-\theta)-\theta D-k_W\big)=\mathbf{-98,1}\\[4pt]
O_7:\quad \N{u_N} &= R=\mathbf{10}
\end{align*}
Sans bonus, ces valeurs étaient $-167,4$ ; $-103$ ; $-103$ ; $0$.

\subsection{La résolution}
\begin{center}
\small
\begin{tabular}{@{}lll@{}}
\toprule
Qui décide & Comparaison & Choix \\
\midrule
\N{Nixon} face au Post en relais & $O_1:-165,69$ contre $O_2:-201$ & poursuit \\
\W{Post} en relais & $O_1:0,56$ contre $O_3:0$ & publie \\
\N{Nixon} face au Times & $O_1:-165,69$ contre $O_4:-200$ & poursuit \\
\N{Nixon} face au Post seul & $O_5:-98,1$ contre $O_6:-200$ & poursuit \\
\W{Post} seul & $O_5:-3,4$ contre $O_7:0$ & se tait \\
\T{Times} & $O_1:6,06$ contre $O_7:0$ & publie \\
\bottomrule
\end{tabular}
\end{center}
\begin{resultat}
\textbf{Résultat B : $O_1$, c'est 1971.} Même résultat que sans bonus.
\end{resultat}

\section{Pourquoi le bonus change si peu}

\subsection{Une réécriture qui explique tout}
On regroupe les termes en $\sigma$ :
\[
R\,(1-\sigma)-D\,\sigma \;=\; R-R\,\sigma-D\,\sigma \;=\; \underbrace{R}_{\text{constante}}\;-\;\underbrace{(D+R)}_{\text{nouveau dommage}}\,\sigma .
\]
Le bonus fait donc deux choses :
\begin{enumerate}[leftmargin=*,itemsep=1pt]
\item \textbf{il ajoute $R=10$ à toutes les feuilles de Nixon}. Cela ne change aucune comparaison : si $a>b$, alors $a+10>b+10$ ;
\item \textbf{il fait passer le dommage de $D$ à $D+R$}, de $200$ à $210$ : un secret gardé compte un peu plus pour Nixon, qui a un peu plus envie de poursuivre.
\end{enumerate}
Vérification sur $O_1$ avec $\theta=0,3$ : avec $D=210$ et sans bonus, $0,3\times(-211)+0,455\times(-212)+0,245\times(-65)=-175,69$ ; on ajoute $10$ : $-165,69$. C'est bien la même valeur.

\subsection{Les seuils de Nixon avec bonus}
Il suffit de remplacer $D$ par $D+R$ dans les seuils de Nixon. Ceux des journaux ne changent pas.

\textbf{Nixon poursuit le Post en relais} si
\[
q<1-\frac{k_W}{(1-\theta)D}\ \ \text{(sans bonus)}\qquad\text{devient}\qquad q<1-\frac{k_W}{(1-\theta)(D+R)}\ \ \text{(avec bonus)}.
\]
Avec $\theta=0,3$ : $q<99,29\,\%$ devient $q<99,32\,\%$.

\textbf{Nixon poursuit le Times} (Post en relais) si
\[
k_T+(1-p)\,k_W<(1-p)(1-q)(1-\theta)\,D\qquad\text{devient}\qquad k_T+(1-p)\,k_W<(1-p)(1-q)(1-\theta)\,(D+R).
\]
Avec $\theta=0,3$ : $1,7<34,3$ devient $1,7<36,0$.

Avec $D=200$, passer à $210$ ne change aucune décision.

\section{Quand le bonus compte : un petit dommage}
Dans le dessin, le dommage de Nixon valait $D=5$. Le bonus pèse alors lourd : $D+R=15$, trois fois plus.

\paragraph{Sans bonus ($D=5$, $\theta=0,3$).} Nixon poursuit-il le Times ?
\[
k_T+(1-p)\,k_W=1+0,7\times1=1,7\quad\text{contre}\quad(1-p)(1-q)(1-\theta)D=0,7\times0,35\times0,7\times5\approx0,86 .
\]
Le procès coûte plus qu'il ne rapporte : \textbf{Nixon laisse faire}, tout sort sans procès ($O_4$).

\paragraph{Avec bonus ($D=5$, $R=10$, $\theta=0,3$).}
\[
1,7\quad\text{contre}\quad 0,7\times0,35\times0,7\times(5+10)\approx2,57 .
\]
Maintenant le procès vaut la peine : \textbf{Nixon poursuit}, et on retrouve 1971 ($O_1$). Son gain en $O_1$ vaut
\[
0,3\times(-6)+0,455\times(-7)+0,245\times(7-1,5-2)=-1,8-3,185+0,8575\approx-4,13>-5=O_4 .
\]
\begin{simple}
Le bonus mesure \emph{combien Nixon tient au secret}. Quand le dommage d'une fuite est déjà énorme ($200$), $10$ de plus ne change rien. Quand il est petit ($5$), le bonus devient décisif.
\end{simple}

\section{Récapitulatif et conseil}
\begin{center}
\small
\begin{tabular}{@{}lcccc@{}}
\toprule
 & \multicolumn{2}{c}{$D=200$} & \multicolumn{2}{c}{$D=5$} \\
 & sans $\theta$ & $\theta=0,3$ & sans $\theta$ & $\theta=0,3$ \\
\midrule
Sans bonus & $O_7$ & $O_1$ & $O_7$ & $O_4$ \\
Avec bonus $R=10$ & $O_7$ & $O_1$ & $O_7$ & $\mathbf{O_1}$ \\
\bottomrule
\end{tabular}
\end{center}
\begin{simple}
\textbf{Conseil.} Si le groupe garde le $+10$, l'écrire comme un bonus $R(1-\sigma)$ dans l'équation de Nixon, avec $\theta$. Phrase possible pour le rapport : « Nixon gagne $R=10$ quand le secret est gardé, en proportion de la part des documents qui reste secrète. » Sinon, laisser $0$ en $O_7$ : avec $D=200$, c'est équivalent et plus simple.
\end{simple}

\end{document}
